\section{\texorpdfstring{\sys}{WitnessRAG}}
\label{sec:method}

\sys treats memory retrieval as the search for facts that jointly satisfy
the information requirements of a question. It operates in two stages
(\cref{fig:overview}). \emph{Memory construction} organizes the
interaction history into a graph of dated propositions linked to their
source passages. \emph{Memory retrieval} grounds a plan in retrieved
evidence, searches the graph for witnesses of that plan, verifies them,
and assembles a compact context for the reader. Throughout, the planner
specifies \emph{what} must be found, and the executor determines
\emph{how} stored facts connect to satisfy it.

\begin{figure*}[t]
  \centering
  \includegraphics[width=\textwidth]{figures/overview-standalone.pdf}
  \caption{\textbf{Overview of \sys.}
  (1) \emph{Memory construction}: dialogue histories are turned into
  dated propositions, stored both as statements and as graph triples
  linked to their source utterances. (2) \emph{Memory retrieval}: a
  planner grounded in retrieved evidence specifies the relational
  requirements of the question, its answer type, and a temporal reference.
  An executor searches the graph for witnesses under consistent bindings;
  a verifier checks candidates against their sources, and gaps or
  rejections are fed back to the planner. Accepted witnesses guide the
  selection of propositions and passage summaries passed to the reader.
  Blue outlines denote LLM operations, green outlines denote retrieval or
  deterministic processing, dotted links denote provenance, and dashed
  arrows denote feedback.}
  \Description{The architecture has two stages: memory construction and memory retrieval. Memory construction
  extracts propositions and indexes their time and importance metadata.
  Source passages and the dated proposition graph are connected by
  provenance. A question triggers evidence retrieval and an LLM planner.
  The planner specifies conjunctive requirements, answer type, temporal
  reference and weights. A graph executor finds consistent bindings without
  LLM calls. A conditional verifier checks source utterances, returning
  gaps or rejections to the planner. Accepted support and relevance-ranked
  candidates guide context selection. The reader receives dated statements
  and passage summaries and generates an answer.}
  \label{fig:overview}
\end{figure*}


\subsection{Memory Construction}
\label{sec:method-register}

The memory pairs a passage store $\passages$ with a fact graph
$\mathcal G=(\mathcal N,\facts)$. Passages preserve the original dialogue
and its session structure. From them, an \ac{LLM} extracts
self-contained propositions, each represented both as a natural-language
statement and as a subject--relation--object triple: the statement is
what the reader consumes, and the triple is what the executor searches.
Entity normalization merges recurring mentions, so that facts recorded
in different sessions can take part in the same derivation.

Each fact is annotated with an event interval, an importance score, and a
link to its source utterance. Temporal expressions are resolved relative
to the session in which they occur, which distinguishes when an event
happened from when it was mentioned. Importance is a
question-independent estimate of the salience expressed in the source
utterance. We write $\mem=(\passages,\facts)$ for the resulting memory
and $\mathrm{src}(W)$ for the source passages of a set of facts $W$.
These provenance links allow any witness to be traced back to the
utterances from which its facts were extracted. Construction is performed
once per history and shared across all questions about it.

\subsection{Evidence-Grounded Planning}
\label{sec:method-plan}

Retrieval starts with a hybrid dense--lexical search, fused by
\ac{RRF}, that yields an initial pool of evidence. The planner, an
\ac{LLM}, receives the question together with dated facts from this pool
and the relations and entities available in memory. Planning
\emph{after} an initial retrieval lets the model phrase its query in the
vocabulary of the stored history rather than that of the question.

A plan $\pi$ consists of a conjunctive query in the form of \cref{eq:cq},
the expected answer type (e.g., an entity, a date, or a set of values),
a temporal reference, and ranking weights. Shared variables encode
dependencies between facts: the value that satisfies one atom constrains
the facts that can satisfy another. Operations such as comparison or
date arithmetic are not executed over the graph; the plan retrieves
their premises, and the reader performs them.

The temporal reference designates the period against which memories are
ranked: the present, the beginning of the history, or an interval named
in the question. Together with the weights, it allows the same memory to
be searched differently for questions about recent events, first
occurrences, or specific periods. Once a plan is available, the evidence
pool is expanded under the corresponding ranking.

\subsection{Reference-Relative Ranking}
\label{sec:method-score}

Passage retrieval and witness search share a scoring function that
combines relevance, importance, and temporal proximity:
\begin{equation}
 s(x\mid q,\pi)=
 \lambda_s\,\mathrm{sim}(x)
 +\lambda_i\,\kappa_x
 +\lambda_t\,g(x;\rho),
 \label{eq:score}
\end{equation}
where $\kappa_x$ is the importance of item $x$, $\rho$ is the temporal
reference of $\pi$, and the nonnegative weights, chosen by the planner,
sum to one. For a passage, $\mathrm{sim}(x)$ is its normalized retrieval
score; for a fact, it measures how well the fact matches the
corresponding query atom.

Temporal proximity decays with the distance between the event interval
$I_x$ and the reference:
\begin{equation}
 g(x;\rho)=\exp\!\left(-\frac{\delta(I_x,\rho)}{h_\rho}\right),
 \label{eq:proximity}
\end{equation}
where $\delta$ vanishes for overlapping intervals and otherwise measures
their separation, and $h_\rho$ sets the decay scale. With the present as
reference, \cref{eq:proximity} reduces to recency; with an earlier
reference, it favors events near that period. Time thus acts as a
question-dependent preference rather than a universal bias toward recent
memories.

\subsection{Witness Search}
\label{sec:method-prove}

The executor searches for witnesses of the planned query without invoking
a language model. As defined in \cref{sec:preliminaries}, a witness is a
set of facts that instantiates every atom under a single consistent
assignment. Matching is semantic rather than exact, which absorbs
differences between the planner's phrasing and the stored propositions;
consistency is strict, so that once a variable is bound, all subsequent
atoms must agree with it. This prevents facts that are individually
plausible from being combined into an incoherent answer.

Search explores partial assignments with a bounded beam, starting from
the most constrained atoms. For a witness $W=\{f_1,\ldots,f_m\}$, where
$f_j$ matches the $j$-th atom, the search minimizes
\begin{equation}
 J(W)=-\sum_{j=1}^{m}\log s(f_j\mid q,\pi)
       +\gamma\,|\mathrm{src}(W)|.
 \label{eq:witness-search}
\end{equation}
The first term favors derivations whose atoms are all well matched, since
a single weak link dominates the sum; the second discourages support
that is dispersed across unnecessarily many passages. A candidate is
accepted only when its matching strength is sufficient on its own, so
that temporal or importance preferences can reorder plausible witnesses
but cannot rescue a poor match. For set-valued answers, support is
assessed separately for each member.

When no consistent extension exists, search reports the unsatisfied atom
together with the bindings established so far. This \emph{gap} states
precisely which requirement remains unresolved and gives the planner a
concrete target for revision.

\subsection{Verification and Replanning}
\label{sec:method-verify}

Witness search establishes that stored facts satisfy the plan, but not
that the plan and the extracted facts faithfully reflect the dialogue.
A verifier therefore checks candidate answers against the source
utterances of their witnesses, assessing whether the evidence concerns
the intended entity, period, and answer type. Members of a set-valued
answer are verified individually, so that a subset can be retained.
Verification is applied selectively, when the support is not already
evident in the retrieved evidence.

If no candidate is accepted, the gaps and the verifier's feedback are
returned to the planner, which may revise a relation, an entity, or the
temporal interpretation before the executor searches again
(\cref{alg:search}). The cycle ends when a witness is accepted or the
cycle budget is exhausted, in which case retrieval falls back to
relevance ranking alone.

\begin{algorithm}[t]
\caption{Witness-guided memory retrieval}
\label{alg:search}
\small
\begin{algorithmic}[1]
\Require Question $q$, memory $\mem$, cycle limit $R$, context budget $b$
\State $E\gets\search(q,\mem)$; $\Delta\gets\emptyset$; $W^+\gets\emptyset$
\For{$c=1,\ldots,R$}
  \State $\pi\gets\plan(q,E,\Delta)$
  \State $E\gets E\cup\search(q,\mem;\pi)$
  \State $(W,\Delta)\gets\prove(\pi,\mem,E)$
  \State $(W^+,\Delta)\gets\verify(q,\pi,W,\Delta)$
  \If{$W^+\neq\emptyset$} \State \textbf{break} \EndIf
\EndFor
\State $\ctx\gets\op{Assemble}(q,\mem,\pi,W^+,b)$
\State \Return $\mathrm{Read}(q,\ctx)$
\end{algorithmic}
\end{algorithm}

\subsection{Context Assembly and Cost}
\label{sec:method-respond}
\label{sec:method-guarantees}

The reader's context combines dated propositions with summaries of their
source passages. Propositions are selected by their membership in
accepted witnesses and by their relevance to the question and the plan,
favoring coverage over redundancy within the budget $b$. They are grouped
by session and annotated with their event dates, while passage summaries,
computed independently of the question and reused across queries,
restore surrounding context that individual propositions omit. Witnesses
thus prioritize evidence without excluding other relevant material. The
reader formulates the answer and carries out any comparison or date
arithmetic the plan has deferred to it.

With at most $R$ cycles, retrieval issues at most $R$ planning and $R$
verification calls, followed by a single reader call; graph execution
requires none. Extraction and summarization are one-time costs amortized
over all questions on the same history. Language models are thus used
where interpretation is required, namely in planning, verification, and
answer generation, while relational composition is carried out by the
executor.
